We presented \textbf{MobileDeepfake}, a two\hyp stage cascaded deepfake detection system engineered specifically for mobile deployment with explicit controls over accuracy, compute, and false negative risk. Our system addresses the fundamental tension between detection performance and on\hyp device constraints by combining a lightweight MobileNetV4 Stage~1 filter with a powerful EfficientNetV2 Stage~2 expert, routing samples between stages based on calibrated confidence thresholds. Through multi\hyp dataset training, hard example mining, temperature scaling calibration, and cost\hyp sensitive threshold tuning, we achieve strong discrimination on combined validation data (Stage~2 AUC~\num{0.9633}, F1~\num{0.8930}) while maintaining deployability on commodity mobile hardware.

\paragraph{Key Findings.}
Our systematic evaluation reveals several important insights:

\begin{itemize}[leftmargin=*,itemsep=2pt,parsep=0pt]
  \item \textbf{Cascade architecture enables accuracy-efficiency trade-offs.} The tuned cascade achieves near\hyp perfect recall (FNR~\num{0.006}) with minimal escalation (\num{1.16}\% of samples routed to Stage~2) under the best threshold configuration (low/high~=~0.05/0.55). This represents a 5.2× improvement in FNR compared to Stage~1\hyp only operation with only 9\% increase in average compute, demonstrating that adaptive inference can reconcile competing deployment constraints.

  \item \textbf{Component contributions are well\hyp understood through ablation studies.} Our comprehensive ablations (Section~6) isolate the impact of hard example mining (+0.76\% F1), temperature scaling calibration (-78.9\% ECE), and meta\hyp model fusion (very small and inconsistent changes relative to the best single expert). These controlled experiments validate each design choice and expose the limits of ensemble fusion, justifying our decision to deploy a simpler two\hyp stage cascade rather than relying on a separate meta\hyp model in the default system.

  \item \textbf{Per\hyp dataset performance varies significantly.} Detailed analysis across CelebDF\hyp v2, FaceForensics++, DFDC, and DeeperForensics\hyp 1.0 reveals that dataset heterogeneity (recording quality, manipulation diversity) strongly influences detection difficulty. Controlled datasets yield higher absolute AUC but can still challenge detectors through manipulation quality, while heterogeneous datasets exhibit lower AUC due to within\hyp dataset variation. Fixed thresholds produce wildly different operating points (FNR ranging from 16\% to 67\% at threshold 0.70), underscoring the need for adaptive threshold tuning.

  \item \textbf{Cross\hyp dataset generalization remains challenging.} On Deepfake\hyp Eval\hyp 2024, our calibrated cascade exhibits substantially degraded performance (F1~$<$0.30, FNR~$\approx$73--79\%) compared to in\hyp distribution validation, aligning with known failure modes on in\hyp the\hyp wild media. This gap stems from preprocessing mismatches, underrepresented generation methods, and dataset\hyp specific artifact learning. Even bypassing Stage~1 (Stage~2\hyp only configuration) does not close the gap, highlighting fundamental generalization challenges that require domain adaptation rather than simple threshold adjustment.

  \item \textbf{Robustness to corruptions shows complex patterns.} Our perturbation sweeps (JPEG, noise, blur, brightness) reveal non\hyp monotonic trends: moderate compression sometimes improves detection by accentuating artifacts, while extreme degradations increase escalation rates as more samples fall into the uncertainty band. These findings motivate per\hyp perturbation calibration and emphasize that reporting Stage~2 escalation rates alongside accuracy metrics is essential for understanding efficiency\hyp robustness trade\hyp offs.

  \item \textbf{Error patterns inform future improvements.} Systematic error analysis identifies recurring failure modes: false positives concentrate in heavily compressed real videos and unusual lighting/makeup, while false negatives stem from high\hyp quality GANs, partial manipulations, out\hyp of\hyp distribution generators, and post\hyp processing. These patterns suggest concrete mitigation strategies (data augmentation, domain adaptation, temporal consistency checks) for future work.

  \item \textbf{GenConViT hybrid architecture underperformed expectations.} Our experiments with GenConViT as an additional Stage~2 expert revealed that hybrid CNN\hyp Transformer architectures, while promising in theory, faced practical challenges including training instability, long training times, and domain\hyp specific biases from pretrained weights under our setup. The simpler EfficientNetV2\hyp B3 baseline matched or exceeded GenConViT performance with significantly better training stability, so our default cascade relies on EfficientNetV2 while GenConViT is retained as an exploratory component (see Appendix for details).
\end{itemize}

\paragraph{Deployment Implications.}
Our exported TorchScript artifacts are compact (Stage~1: \SI{37.3}{MB}, Stage~2: \SI{49.6}{MB}; Table~\ref{tab:mobile_artifacts_auto}) and suitable for Android integration via PyTorch Mobile. The system exposes explicit operational dials---Stage~1 leakage rate $\pi_{\mathrm{leak}}$, Stage~2 escalation rate $\pi_2$, and tunable thresholds $(\tau_{low}, \tau_{high})$---that allow operators to configure deployment scenarios ranging from safety\hyp first moderation (low FNR, moderate escalation) to bandwidth\hyp constrained offline tools (higher FNR, minimal escalation). Unlike opaque single\hyp model detectors, our cascade makes risk/efficiency trade\hyp offs measurable and auditable, facilitating deployment decisions aligned with specific platform requirements.

The full pipeline---from raw dataset preprocessing through multi\hyp dataset training, hard example mining, calibration, threshold tuning, robustness evaluation, and mobile export---is reproducible from provided scripts and manifests. This end\hyp to\hyp end auditability enables practitioners to adapt the system to new datasets, tune operating points for domain\hyp specific requirements, and integrate detection into production systems with clear understanding of expected behavior and failure modes.

\paragraph{Limitations and Future Work.}
Several limitations warrant further investigation:

\begin{itemize}[leftmargin=*,itemsep=2pt,parsep=0pt]
  \item \textbf{Distribution shift and adaptation.} While our multi\hyp dataset training improves generalization compared to single\hyp dataset baselines, Deepfake\hyp Eval\hyp 2024 results demonstrate that significant performance gaps remain on truly in\hyp the\hyp wild media. Future work should explore lightweight domain adaptation techniques (e.g., domain\hyp adaptive batch normalization~\cite{li2021dabn}, unsupervised feature alignment) that can improve cross\hyp dataset robustness without inflating on\hyp device compute costs. Per\hyp domain calibration and dynamic threshold adjustment based on automatic domain detection represent promising directions.

  \item \textbf{Temporal consistency and video\hyp level aggregation.} Our frame\hyp level classification approach does not exploit temporal consistency or cross\hyp frame dependencies. Incorporating temporal modeling (e.g., LSTM aggregation over frame\hyp level predictions, 3D CNNs for spatiotemporal features) could improve robustness to partial manipulations and reduce false positives from transient artifacts. The cascade architecture naturally supports such extensions: Stage~1 can pre\hyp filter frames, and Stage~2 can refine segment boundaries under a compute budget.

  \item \textbf{Model compression and latency optimization.} While our quantized exports reduce model size, further compression through knowledge distillation, neural architecture search, or pruning could shrink Stage~2 latency without sacrificing accuracy. Exploring lightweight transformer variants or efficient attention mechanisms might unlock the benefits of hybrid architectures (like GenConViT) while addressing the training instability and compute overhead we encountered.

  \item \textbf{On\hyp device latency measurements.} Our current evaluation reports FLOPs and theoretical latency estimates but lacks comprehensive on\hyp device benchmarks across diverse mobile hardware (low\hyp end vs.\ flagship devices, ARM vs.\ x86 architectures, with/without NNAPI acceleration). Conducting such benchmarks would validate real\hyp world deployment feasibility and identify hardware\hyp specific bottlenecks.

  \item \textbf{Fairness and bias auditing.} We do not stratify performance metrics by demographic groups, languages, or content domains. Future work should incorporate fairness diagnostics (group\hyp disaggregated FNR/FPR, subgroup robustness) and, where appropriate, apply bias mitigation techniques to reduce disparities. This is particularly important for deployments affecting diverse user populations.

  \item \textbf{Adversarial robustness.} Our robustness evaluations cover natural corruptions (compression, noise, blur) but not targeted adversarial attacks (e.g., C\&W, PGD perturbations). Evaluating and hardening the cascade against adaptive adversaries---especially those aware of the two\hyp threshold routing mechanism---remains critical for high\hyp security deployments.

  \item \textbf{Longitudinal monitoring and model updates.} As generative models evolve, detection performance may degrade over time. Establishing protocols for continuous monitoring (tracking escalation rates, FNR incidents, user feedback), periodic retraining on new fake samples, and seamless model updates without disrupting deployed systems represents an important operational challenge.
\end{itemize}

\paragraph{Broader Impact.}
Beyond technical contributions, \textbf{MobileDeepfake} demonstrates that research prototypes can be engineered into deployable, auditable systems without sacrificing methodological rigor. By coupling cascade design with explicit threshold controls, calibration, systematic ablations, and reproducible scripts, we bridge the gap between academic deepfake detection research and practical mobile deployment. Our open\hyp source pipeline (code, manifests, scripts available at repository URL) enables researchers and practitioners to build upon this foundation, adapt it to new domains, and contribute improvements back to the community.

The shift toward on\hyp device processing aligns with growing privacy expectations and regulatory requirements (e.g., GDPR, data localization laws) that limit server\hyp side media analysis. Our system shows that strong deepfake detection is achievable within these constraints, empowering users with local verification tools rather than requiring trust in opaque cloud services. This democratization of detection capabilities---putting effective tools directly in users' hands---represents a meaningful step toward resilient, privacy\hyp preserving defenses against synthetic media manipulation.

\paragraph{Conclusion.}
Deepfakes will continue to evolve in sophistication and accessibility, necessitating equally adaptive and deployable countermeasures. \textbf{MobileDeepfake} provides a practical methodology for building such systems: multi\hyp dataset training for generalization, cascade routing for efficiency, explicit threshold tuning for risk control, systematic evaluation including OOD benchmarks and robustness sweeps, and reproducible mobile export. While challenges remain---especially cross\hyp dataset gaps and adversarial robustness---our work establishes a foundation for principled, auditable mobile deepfake detection that balances accuracy, efficiency, and operational transparency. We invite the community to build upon this system, contribute improvements, and collectively advance the state of deployable detection in the face of rapidly advancing generative models.
